% Gerado por scripts/etapa3_xgboost.py — requer \usepackage{booktabs}
\begin{table}[htbp]
\centering
\small
\caption{XGBoost — quantis empíricos versus normais (sensibilidade: estoque arredondado para cima)}
\label{tab:xgb_empirico_vs_normal_sensibilidade}
\begin{tabular}{lrrrrrrr}
\toprule
Nível (razão) & Custo empírico & Custo normal & Diferença [IC 95\%] & Sobra empírico & Sobra normal & Falta empírico & Falta normal \\
\midrule
80\% (1:4) & 2,727 & 2,847 & +4,4 [+3,9; +4,9] & 1,740 & 2,090 & 0,247 & 0,189 \\
90\% (1:9) & 3,838 & 3,836 & -0,1 [-0,4; +0,3] & 2,648 & 2,758 & 0,132 & 0,120 \\
95\% (1:19) & 4,985 & 4,918 & -1,3 [-1,6; -1,1] & 3,612 & 3,330 & 0,072 & 0,084 \\
\bottomrule
\end{tabular}
\par\smallskip
{\footnotesize Nota: teste de 90 dias (26/01 a 24/04/2016), 5.748 séries. Diferença: custo normal em relação ao empírico, em \% (positivo = normal mais caro), com intervalo de 95\% de confiança entre colchetes, item como unidade (1.437 itens). Sobra e falta em unidades por série-dia. custo por série-dia com custo de sobra = 1; razão 1:4, 1:9 e 1:19 (sobra:falta) corresponde aos níveis de serviço de 80\%, 90\% e 95\%. Estoque = quantil arredondado para cima.}
\end{table}
